%================================
% note-setup-printer.tex
% Printer-friendly version
%================================

\usepackage{amsmath,amsthm,amssymb}
\usepackage{mathtools}
\usepackage{mathrsfs}
\usepackage{bm}
\usepackage{extarrows}

\usepackage[a4paper, margin=1cm, footskip=0.5cm]{geometry}
\usepackage{float}
\usepackage{indentfirst}
\usepackage{booktabs,multirow}
\usepackage[shortlabels,inline]{enumitem}
\usepackage{appendix}

%================================
% FIX GLOBALE PER LISTE DOPO I TEOREMI
%================================

\makeatletter
\AddToHook{env/itemize/begin}{%
    \if@inlabel\mbox{}\par\vspace{-\baselineskip}\fi
}
\AddToHook{env/enumerate/begin}{%
    \if@inlabel\mbox{}\par\vspace{-\baselineskip}\fi
}
\makeatother

%================================
% PACCHETTI GRAFICI
%================================

\usepackage[dvipsnames]{xcolor}
\usepackage{pifont}
\usepackage{graphicx}

\graphicspath{
    {./figure/}{./figures/}
    {./image/}{./images/}
    {./graphic/}{./graphics/}
    {./picture/}{./pictures/}
}

\usepackage{subcaption}

%================================
% ALGORITMI E CODICE
%================================

\usepackage[ruled,linesnumbered,noline]{algorithm2e}
\usepackage{listings}

\lstdefinestyle{simpleStyle}{
    basicstyle=\ttfamily\small,
    breaklines=true,
    keywordstyle=\color{blue},
    identifierstyle=\color{black},
    stringstyle=\color{violet},
    commentstyle=\color[RGB]{34,139,34},
    showstringspaces=false,
    numbers=left,
    numbersep=2em,
    numberstyle=\footnotesize,
    frame=single,
    framesep=1em,
}

\lstset{style=simpleStyle}

%================================
% HYPERREF
%================================

\usepackage{hyperref}

\hypersetup{
    colorlinks=true,
    linkcolor=black,
    urlcolor=cyan,
    citecolor=black
}

%================================
% TEOREMI E TCOLORBOX
%================================

\usepackage[most]{tcolorbox}
\usepackage{keytheorems}

%================================
% STILE GENERALE: SOLO BARRA SINISTRA
%================================

\tcbset{
    theorem box/.style={
        enhanced,
        breakable,
        sharp corners,
        boxrule=0pt,
        leftrule=2pt,
        rightrule=0pt,
        toprule=0pt,
        bottomrule=0pt,
        colback=white,
        boxsep=3pt,
        left=5pt,
        right=2pt,
        top=3pt,
        bottom=3pt,
        before skip=6pt,
        after skip=6pt,
    },
}

%================================
% STILI UFFICIALI KEYTHEOREMS
%================================

\newkeytheoremstyle{breakplain}{
    inherit-style=plain,
    headpunct={},
    break
}

\newkeytheoremstyle{breakdefinition}{
    inherit-style=definition,
    headpunct={},
    break
}

\newkeytheoremstyle{breakremark}{
    inherit-style=remark,
    headpunct={},
    break
}

\newkeytheoremstyle{notestyle}{
    headfont=\color{orange!80!black}\bfseries,
    bodyfont=\normalfont,
    spaceabove=3pt,
    spacebelow=3pt,
    headpunct={},
    break,
    qed=\ensuremath{\color{orange!80!black}\diamond}
}

\newkeytheoremstyle{solutionstyle}{
    headfont=\bfseries,
    bodyfont=\normalfont,
    spaceabove=3pt,
    spacebelow=3pt,
    headpunct={},
    break,
    qed=\ensuremath{\square}
}

%================================
% DEFINIZIONE AMBIENTI
%================================

% TEOREMI - Viola

\newkeytheorem{theorem}[
    style=breakplain,
    name=Teorema,
    numbered=true,
    parent=section,
    tcolorbox-no-titlebar={
        theorem box,
        colframe=RoyalPurple
    }
]

\newkeytheorem{theorem*}[
    style=breakplain,
    name=Teorema,
    numbered=false,
    tcolorbox-no-titlebar={
        theorem box,
        colframe=RoyalPurple
    }
]

% PROPOSIZIONI - Viola

\newkeytheorem{prop}[
    style=breakplain,
    name=Proposizione,
    numbered=true,
    sibling=theorem,
    tcolorbox-no-titlebar={
        theorem box,
        colframe=RoyalPurple
    }
]

\newkeytheorem{prop*}[
    style=breakplain,
    name=Proposizione,
    numbered=false,
    tcolorbox-no-titlebar={
        theorem box,
        colframe=RoyalPurple
    }
]

% COROLLARI - Blu

\newkeytheorem{corollary}[
    style=breakplain,
    name=Corollario,
    numbered=true,
    sibling=theorem,
    tcolorbox-no-titlebar={
        theorem box,
        colframe=NavyBlue
    }
]

\newkeytheorem{corollary*}[
    style=breakplain,
    name=Corollario,
    numbered=false,
    tcolorbox-no-titlebar={
        theorem box,
        colframe=NavyBlue
    }
]

% LEMMI - Blu

\newkeytheorem{lemma}[
    style=breakplain,
    name=Lemma,
    numbered=true,
    sibling=theorem,
    tcolorbox-no-titlebar={
        theorem box,
        colframe=NavyBlue
    }
]

\newkeytheorem{lemma*}[
    style=breakplain,
    name=Lemma,
    numbered=false,
    tcolorbox-no-titlebar={
        theorem box,
        colframe=NavyBlue
    }
]

% CLAIM - Blu

\newkeytheorem{claim}[
    style=breakplain,
    name=Claim,
    numbered=true,
    sibling=theorem,
    tcolorbox-no-titlebar={
        theorem box,
        colframe=NavyBlue
    }
]

\newkeytheorem{claim*}[
    style=breakplain,
    name=Claim,
    numbered=false,
    tcolorbox-no-titlebar={
        theorem box,
        colframe=NavyBlue
    }
]

% DEFINIZIONI - Verde

\newkeytheorem{definition}[
    style=breakdefinition,
    name=Definizione,
    numbered=true,
    parent=section,
    tcolorbox-no-titlebar={
        theorem box,
        colframe=ForestGreen
    }
]

\newkeytheorem{definition*}[
    style=breakdefinition,
    name=Definizione,
    numbered=false,
    tcolorbox-no-titlebar={
        theorem box,
        colframe=ForestGreen
    }
]

% ESEMPI - Terra

\newkeytheorem{example}[
    style=breakdefinition,
    name=Esempio,
    numbered=true,
    parent=section,
    tcolorbox-no-titlebar={
        theorem box,
        colframe=RawSienna
    }
]

\newkeytheorem{example*}[
    style=breakdefinition,
    name=Esempio,
    numbered=false,
    tcolorbox-no-titlebar={
        theorem box,
        colframe=RawSienna
    }
]

% ESERCIZI - Rosso

\newkeytheorem{esercizio}[
    style=breakdefinition,
    name=Esercizio,
    numbered=true,
    parent=section,
    tcolorbox-no-titlebar={
        theorem box,
        colframe=WildStrawberry!70
    }
]

\newkeytheorem{esercizio*}[
    style=breakdefinition,
    name=Esercizio,
    numbered=false,
    tcolorbox-no-titlebar={
        theorem box,
        colframe=WildStrawberry!70
    }
]

% OSSERVAZIONI - Senza riquadro

\newkeytheorem{oss}[
    style=breakremark,
    name=Osservazione,
    numbered=true,
    parent=section
]

\newkeytheorem{oss*}[
    style=breakremark,
    name=Osservazione,
    numbered=false
]

% NOTE - Arancione

\newcommand{\notehandmark}{%
    \raisebox{0.05ex}{\large\ding{45}}\hspace{0.4em}%
}

\newkeytheorem{note}[
    style=notestyle,
    name={\notehandmark Nota},
    numbered=true,
    parent=section,
    tcolorbox-no-titlebar={
        theorem box,
        colframe=orange!80!black
    }
]

\newkeytheorem{note*}[
    style=notestyle,
    name={\notehandmark Nota},
    numbered=false,
    tcolorbox-no-titlebar={
        theorem box,
        colframe=orange!80!black
    }
]

% SOLUZIONI - Nero

\newkeytheorem{solution}[
    style=solutionstyle,
    name=Soluzione,
    numbered=true,
    parent=section,
    tcolorbox-no-titlebar={
        theorem box,
        colframe=black!65
    }
]

\newkeytheorem{solution*}[
    style=solutionstyle,
    name=Soluzione,
    numbered=false,
    tcolorbox-no-titlebar={
        theorem box,
        colframe=black!65
    }
]

%================================
% COVER (INVARIATA)
%================================

\usepackage{titling}

\newcommand{\extrainfocontent}{}

\newcommand{\extrainfo}[1]{%
    \renewcommand{\extrainfocontent}{#1}%
}

\newcommand{\makecover}[1]{%
    \begingroup
    \hypersetup{pageanchor=false}
    \begin{titlepage}
        \newgeometry{margin=0in}
        \parindent=0pt

        \includegraphics[width=\linewidth]{#1}

        \vfill

        \begin{center}
            \parbox{0.618\textwidth}{
                \raggedleft
                {\bfseries \Huge \thetitle}
                \\[0.6pt]
                \rule{0.618\textwidth}{4pt}
                \\
            }
        \end{center}

        \vfill

        \begin{center}
            \parbox{0.618\textwidth}{
                \raggedleft\Large
                \begin{tabular}{r}
                    \theauthor \\
                    \thedate
                \end{tabular}
            }
        \end{center}

        \vfill

        \begin{center}
            \parbox[t]{0.7\textwidth}{
                \centering
                \itshape
                \extrainfocontent
            }
        \end{center}

        \vfill
    \end{titlepage}

    \restoregeometry
    \endgroup
}

%================================
% DIMOSTRAZIONI
%================================

\renewcommand{\proofname}{Dimostrazione}

\makeatletter

\renewenvironment{proof}[1][\proofname]{%
    \par
    \pushQED{\qed}
    \normalfont
    \topsep6\p@\@plus6\p@\relax
    \trivlist
    \item[\hskip\labelsep\bfseries #1]
    \ignorespaces\leavevmode\par\nobreak
}{%
    \popQED
    \endtrivlist
    \@endpefalse
}

\makeatother

% Dimostrazioni: barra verde petrolio

\tcolorboxenvironment{proof}{
    theorem box,
    colframe=TealBlue
}

%================================
% PENCIL NOTE
%================================

\newcommand{\pencilnote}[1]{%
    {\color{gray!80!black}#1}%
}

%================================
% FINE TEMPLATE
%================================